\section{Reduction through sourced constraints}\label{d:reduction}

In this section, $g,k$ denote the \emph{polarized} data of
Proposition~\ref{a:polarization}, and $E_g$ is their ADM energy.  Write
$\cD_X=(D_1,D_2,D_3)$ for the three covector rows of the potential
expression of the axial construction; in particular, the
third row contains the Pfaff form $\dd z+\psi\dd\chi-\chi\dd\psi$.
The argument also applies to the following larger class of data, which
will be needed for the equality variations.  The proposition
quantifies over all admissible potential maps with the specified
axis constants, including their compactly supported variations.
Within each individual application the chosen potentials remain
fixed throughout the metric preparation and deformation.  The
constraint densities are always computed from the current $g,k$.

\begin{proposition}[Comparison for sourced data]\label{d:generalized}
Suppose a smooth polarized rotation exterior has one asymptotically
flat end, a compact weakly future trapped boundary, and the topology
and potential regularity of the axial construction.  Assume
$g-\delta=O_2(r^{-1})$, $k=O_1(r^{-3})$,
$|s|+|\cD_X|_g=O(r^{-2})$, $R_g=O(r^{-3-\delta})$ for some
$0<\delta<1$, and finite ADM energy.  The potentials have the axis
and end regularity in Lemma~\ref{a:axis-regularity}.  On each signed
axis tube, $s_1,s_2$ are smooth odd functions of order $y$, and $s_3$
is smooth even of order $y^2$, with smooth longitudinal derivatives.
For each row, the functions $|D_i|_g^2,s_i^2$ and the covector $s_iD_i$
extend smoothly to the three-dimensional exterior through the axis.
The source expressions below use these extensions at axis points;
neither an individual normalized row nor its ordinary norm is
required to extend smoothly.  Suppose, at every point and for every
$w\in T\Omega$ with $|w|_g\le1$, that
\begin{equation}\label{d:sourced}
 8\pi\bigl(\mu+J_{\rm con}(w)\bigr)
 \ge |\cD_X|_g^2+|s|^2+2s\cdot\cD_X(w).
\end{equation}
Let $L>0$ be a lower bound for all crossing lengths in the meridional
area metric of $g$.  For each admissible subextremal model radius $R$,
with $F_*(R)$ and $b=b(R)>0$ as in Lemma~\ref{a:model-parameters},
\begin{equation}\label{d:sourced-comparison}
 E_g\ge F_*(R)+\frac b2\log\frac{L}{2R^2}.
\end{equation}
\end{proposition}

In particular, $J_{\rm con}$ is retained in
Equation~\eqref{d:sourced}; the local momentum density need
not vanish.  The three source rows are estimated separately.  The
neutral geometric deformation needed to prove the proposition is the
following theorem.  Its proof occupies Section~\ref{d:proof}, including
the existence argument and end-flux estimate.

\begin{theorem}[Neutral deformation with a prescribed loss]\label{d:deformation}
Let $(\Omega,g,k)$ be a smooth polarized one-ended exterior, complete
with its compact nonempty boundary included.  Assume the boundary is
strictly future trapped, $k$ has compact support, $g-\delta=O_j(r^{-1})$
for every finite $j$, $R_g=O(r^{-3-\delta})$ for some $0<\delta<1$,
and $g$ has finite ADM energy.  Let $\mathfrak m>0$ be a smooth
invariant function such that
\[
 \mathfrak m\le8\pi(\mu-|J_{\rm con}|_g),\qquad
 \inf_C\mathfrak m>0\quad\text{for every compact }C\subset\overline\Omega,
 \qquad \mathfrak m\ge c r^{-3-\delta}\quad\text{on the end}
\]
for some $c>0$.  There is an invariant reduced exterior
$\Omega_{\rm red}\subset\Omega$, obtained by removing full trapped
regions and retaining the component toward infinity, with smooth
nonempty compact boundary, for which the following holds.  For every
sufficiently small fixed $\eps>0$, and all sufficiently large integers
$n$ along a sequence, there are smooth invariant functions $f_n,t_n$
and positive smooth profiles $l_n(t)$ such that
\begin{gather}\label{d:deformed-metrics}
 \bar g_n=g+l_n(t_n)^2\dd f_n^2,\qquad
 \hat g_n=e^{4t_n}\bar g_n,\qquad
 w_n=\frac{l_n(t_n)\nabla f_n}{\sqrt{1+l_n(t_n)^2|\dd f_n|_g^2}},
 \qquad t_n\ge-\eps,\\
 \frac12e^{4t_n}R_{\hat g_n}
 \ge8\pi\bigl(\mu+J_{\rm con}(w_n)\bigr)-\frac{\mathfrak m}{2}.
 \label{d:deformation-curvature}
\end{gather}
The inner boundary has strictly negative mean curvature in $\hat g_n$
for its normal toward infinity.  The metric $\hat g_n$ is complete
with boundary, is asymptotically flat of order one with at least two
controlled derivatives, and has finite energy satisfying
\begin{equation}\label{d:deformation-energy}
 E_{\hat g_n}\le E_g+o_n(1).
\end{equation}
Here the reduced domain and its collars are fixed before $\eps$ and
$n$.  For each fixed $n$ the functions are obtained by exhaustion of
finite outer truncations; the error in
Equation~\eqref{d:deformation-energy} tends to zero as $n\to\infty$
after that exhaustion, with the prepared data and $\eps$ fixed.
\end{theorem}

\begin{proof}[Proof of Proposition~\ref{d:generalized}, using Theorem~\ref{d:deformation}]
We first construct a strictifying direction, including the required
elliptic solve.  There is a smooth positive invariant function $v$ with
\begin{equation}\label{d:strictifier}
 \nu v=-1\quad\text{on }\partial\Omega,\qquad
 -\Delta_gv-|k|_g|\dd v|_g\ge c\langle r\rangle^{-3-\delta},\qquad
 v=O_2(r^{-1}),
\end{equation}
whose Laplacian is integrable and whose end flux has a finite limit.
Here $\nu$ points into the exterior and $\langle r\rangle$ is a smooth
positive extension of $r$ over the compact part.  Choose a smooth
invariant majorant $d_0\ge|k|_g$ with $d_0=O(r^{-2})$, and positive
smooth weights $\zeta_0,w_0$ of orders $-2,-3-\delta$, respectively,
with the corresponding symbol bounds.  Solve
\begin{equation}\label{d:strictifier-equation}
 -\Delta_gv=d_0\sqrt{|\dd v|_g^2+\zeta_0^2}+w_0,\qquad
 \nu v=-1,\qquad v|_{S_R}=0
\end{equation}
on a finite truncation.  Continuation from $d_0=0$ has an invertible
linearization: it is a uniformly elliptic operator with bounded drift,
homogeneous Neumann data on the inner boundary and homogeneous
Dirichlet data on $S_R$.  The maximum principle gives uniqueness and
nonnegativity.  On each fixed truncation, $W^{2,p}$ estimates and
interpolation bound the solution by its supremum.  If these suprema
were unbounded during continuation, division by them and compactness
would give a nonzero homogeneous bounded-drift solution with
homogeneous mixed data, contradicting the maximum principle.  This
proves existence on each truncation.

The bounds persist under exhaustion.  On a fixed far region, a
sufficiently large multiple of
$r^{-1}(1-Cr^{-\delta})$, multiplied by one plus the compact-core
supremum, is an upper barrier: the favorable term in
$\Delta(r^{-1-\delta})$ dominates the metric and drift errors of order
$r^{-4}$.  If the core suprema along an exhaustion were unbounded,
normalize once more.  The barrier makes the limit tend to zero at
infinity, while its value is one somewhere on the core.  Hence it
attains a positive global maximum.  It satisfies a homogeneous
bounded-drift equation with homogeneous inner Neumann data, so the
strong maximum principle and
boundary point lemma exclude this.  Exhaustion therefore gives a
positive solution of Equation~\eqref{d:strictifier-equation} with the
outer condition replaced by decay to zero.

For completeness, rescale on an end annulus by $x\mapsto Rx$ and
$v\mapsto Rv(Rx)$.  The value bound just proved gives uniform
$W^{2,p}$ bounds on smaller annuli for $p>3$: the rescaled equation has
bounded coefficients and at most linear growth in the gradient.
The resulting $C^{1,\alpha}$ control makes the right side uniformly
H\"older, so Schauder estimates give the stated $O_2(r^{-1})$ bound.
They also give higher symbol estimates through the orders allowed by
the coefficients.  The equation now gives an integrable Laplacian;
the Euclidean and metric Laplacians differ by integrable errors.
The divergence theorem yields a finite end flux.  These facts prove
Equation~\eqref{d:strictifier}.

For a small constant $\delta_{\rm p}>0$, make the change
\begin{equation}\label{d:strictification}
 g\longmapsto e^{4\delta_{\rm p}v}g,\qquad
 k\longmapsto e^{2\delta_{\rm p}v}k,\qquad
 s_i\longmapsto e^{-2(1+h_i)\delta_{\rm p}v}s_i,\qquad
 (h_1,h_2,h_3)=(1,1,2),
\end{equation}
and keep the potentials fixed.  Identify the old and new unit balls
by $w_{\rm new}=e^{-2\delta_{\rm p}v}w_{\rm old}$.  After multiplying
the new constraint inequality by $e^{4\delta_{\rm p}v}$, its $i$th
source row is its old value times $e^{-4h_i\delta_{\rm p}v}\le1$.
Each row is nonnegative on the unit ball, since
\[
 |D_i|_g^2+s_i^2+2s_iD_i(w)
 =|D_i|_g^2-D_i(w)^2+(s_i+D_i(w))^2\ge0.
\]
This calculation is made off the axis; its left side is a smooth
function on the unit-ball bundle and the inequality extends by
continuity.  Smooth invariant conformal factors preserve the stated
polar orders and all these quadratic source combinations.
The conformal scalar-curvature and momentum-divergence formulas add
to the left side the quantity
\[
 -4\delta_{\rm p}\Delta v-4\delta_{\rm p}^2|\dd v|^2
       +4\delta_{\rm p}k(\nabla v,w_{\rm old}).
\]
Equation~\eqref{d:strictifier} makes this positive, uniformly on compact
sets and at least $c_{\delta_{\rm p}}r^{-3-\delta}$ on the end, after
decreasing $\delta_{\rm p}$ if necessary.  Inside its positive scaling
factor the boundary expansion decreases by $4\delta_{\rm p}$.
Thus the new boundary is strictly future trapped and
Equation~\eqref{d:sourced} has a positive residual.  Crossing
lengths vary by a factor tending to one, scalar-curvature decay is
preserved, and the energy tends to $E_g$ as $\delta_{\rm p}\downarrow0$
by the finite flux of $v$.

Keep $\delta_{\rm p}>0$ fixed.  Cutting off $k$ sufficiently far out
changes the constraints by $O(r^{-4})$, which is smaller than the
strict residual because $\delta<1$.  If the metric does not yet have
all end symbol derivatives, smooth only its far end by averaging
normalized pullbacks under small constant rotations and dilations.
Use smooth kernels commuting with the axial action and the
polarization reflection; for example, the rotation kernel may be
conjugation invariant.  This is smoothing in angle and logarithmic
radius and gives symbol bounds of every order.  The Euclidean linear
part of scalar curvature averages with positive dilation factors;
the nonlinear errors are $O(r^{-4})$ using the original two derivative
bounds.  Hence its positive $r^{-3-\delta}$ lower bound and its
$O(r^{-3-\delta})$ upper order persist.  The source rows are
$O(r^{-4})$, so the full residual persists with a smaller constant.
Patch to the original metric on a fixed distant annulus, choosing
the averaging aperture small enough that the compact patch errors
fit within the residual.  Uniform equivalence tends to identity as
the aperture decreases.  Averaging the linear sphere flux, using
rotation invariance and the dilation scaling, also gives convergence
of the ADM energies.  All preparations preserve polarization.

Choose a smooth positive $\mathfrak m$ no larger than this sourced
residual for any $|w|_g\le1$, with the compact and end lower bounds in
Theorem~\ref{d:deformation}.  Indeed, the residual has the form
$a(x)+b_x(w)$ with $a$ a smooth invariant scalar and $b$ a smooth
invariant covector.  Its minimum on the unit ball is the continuous
positive function $a-|b|_g$, bounded below on compact sets and by
$c r^{-3-\delta}$ on the end.  A smooth positive minorant is obtained
by a partition of unity on the compact part, with a sufficiently
small multiple of $r^{-3-\delta}$ on the end, followed by averaging
over the circle.  No derivative of a row norm is used.  The row
nonnegativity implies also
$\mathfrak m\le8\pi(\mu-|J_{\rm con}|_g)$.  The theorem therefore
applies to the prepared data.  Set
\[
 g_r=e^{4\eps}\hat g_n,\qquad X_r=e^{2(t_n+\eps)}X\ge X,
\]
and renormalize the asymptotic coordinates by the constant scale.
The identity $\bar g_n^{-1}=g^{-1}-w_n\otimes w_n$ gives, row by row,
\begin{equation}\label{d:row-completion}
 |D_i|_g^2+s_i^2+2s_iD_i(w_n)
 =|D_i|_{\bar g_n}^2+(s_i+D_i(w_n))^2.
\end{equation}
Here the calculation again holds off the axis.  Since $f_n$ is a
smooth invariant function, $\bar g_n$ has the same smooth polar
parities as $g$.  The contractions $D_1(w_n),D_2(w_n)$ are smooth
odd of order $y$, and $D_3(w_n)$ is smooth even of order $y^2$.
Thus the completed squares and the new squared row norms extend
smoothly across the axis.  In particular the resulting
scalar-curvature inequality extends there by continuity.
Combining the sourced residual with
Equation~\eqref{d:deformation-curvature} leaves the strictly positive
quantity $\mathfrak m/2$.  Each row of $\cD_{X_r}$ has in addition the
factor $e^{-2h_i(t_n+\eps)}\le1$.  Consequently
\[
 \frac12R_{g_r}>|\cD_{X_r}|_{g_r}^2.
\]
Writing the polarized metric as $g=h+X^2\dd\phi^2$, its meridional
area metric satisfies
\begin{equation}\label{d:length-monotonicity}
 X_r^2h_r=e^{8(t_n+\eps)}X^2(h+l_n^2\dd f_n^2)\ge X^2h.
\end{equation}
Thus every new crossing arc has the old lower length bound, and the
potential and axis regularity required by
Proposition~\ref{r:comparison} are retained.  The end energy obeys
$E_{g_r}\le e^{2\eps}(E_g+o_n(1))$.  Apply that proposition to each
fixed admissible $R$, first exhaust the truncation at fixed $n$, then
let $n\to\infty$, then $\eps\downarrow0$.  Finally undo the far-end
smoothing and cutoff and let $\delta_{\rm p}\downarrow0$.  The
energies and crossing lower bounds converge as just proved, giving
Equation~\eqref{d:sourced-comparison}.
\end{proof}

\begin{corollary}[Numerical inequality]\label{d:numerical}
For the data of Theorem~\ref{m:main},
\[
 m^2\ge\frac{A}{16\pi}+\frac{Q^2}{2}
              +\frac{\pi(Q^4+4J^2)}{A}.
\]
More generally, in the sourced class of
Proposition~\ref{d:generalized}, if $L/2>\sqrt{Q^4+4J^2}$, then
$E_g\ge F_*(\sqrt{L/2})$.
\end{corollary}
\begin{proof}
For the original data, Equation~\eqref{a:constraints} supplies
Equation~\eqref{d:sourced}.  Its scalar equation, with the axial
connection-curvature contribution retained in $s$, gives
$R_g=O(r^{-4})$ from $|\cD_X|+|s|=O(r^{-2})$ and
$k=O_1(r^{-3})$.  This meets the required scalar decay for every
$0<\delta<1$.  The available two metric derivatives and one tensor
derivative suffice even if forming polar components loses a
derivative.  The crossing-area comparison and
the preservation of invariant areas under polarization give
$L=A/(2\pi)$, while $E_g=m$.  If
$A/(4\pi)>\sqrt{Q^4+4J^2}$, choose $R^2=L/2$ in
Equation~\eqref{d:sourced-comparison}; its logarithmic term vanishes.
At the closed endpoint use subextremal radii $R^2\downarrow L/2$;
then $b(R)\to0$ and the same conclusion follows by continuity.  The
formula for $F_*$ in Lemma~\ref{a:model-parameters} and $m>0$ give the
displayed inequality after squaring.  The same choice of $R$ proves
the last assertion for sourced data.  No classification of equality
at the closed endpoint is used here.
\end{proof}
